% ============================================================
% Module 00, Unit 03 — Single-Variable Review & Notation Setup
% Exercises
% Threat Surfaces: Multivariable Calculus for AI Security
% fischer³ Education
% ============================================================

\newcommand{\coursemodule}{Module 00, Unit 03}
\newcommand{\coursetitle}{Single-Variable Review \& Notation Setup}
\newcommand{\coursedate}{February 2026}

\documentclass[11pt, letterpaper]{article}
\input{../../../assets/latex-preamble.tex}

% ============================================================
\begin{document}
% ============================================================

\maketitleblock

% ------------------------------------------------------------
\begin{tcolorbox}[
  colback=ts-light,
  colframe=ts-gray,
  title=Before You Begin,
  fonttitle=\bfseries
]
Complete \texttt{notes.md} for this unit before working these exercises.
If the chain rule or logarithm rules feel shaky, visit the refreshers first:
\texttt{refreshers/single-variable-chain-rule/} and
\texttt{refreshers/exponential-and-log-rules/}.

\medskip
\textbf{Estimated time:} 20--30 minutes \quad
\textbf{Problems:} 4 \quad
\textbf{Difficulty:}
$\star$ Foundational \enspace
$\star\star$ Applied \enspace
$\star\star\star$ Extension
\end{tcolorbox}

\vspace{16pt}


% ============================================================
% PROBLEM 1 — Differentiation
% ============================================================
\begin{problembox}[title={Problem 1 \hfill $\star$}]
Differentiate each function with respect to $x$. Show all intermediate steps,
clearly identifying which rule (chain rule, product rule, or both) you are
applying at each stage.
\end{problembox}

\begin{enumerate}
  \item $f(x) = e^{-3x^2 + 2x}$

  \item $g(x) = x^3 \ln(x^2 + 1)$

  \item $h(x) = \ln\!\Bigl(\dfrac{e^x}{1 + e^x}\Bigr)$
  \quad \textit{Hint: simplify using logarithm rules before differentiating.}

  \item $k(x) = \sqrt{2\pi\sigma^2}\,\exp\!\left(-\dfrac{(x - \mu)^2}{2\sigma^2}\right)$
  \quad \textit{Treat $\mu$ and $\sigma^2$ as constants; differentiate with respect to $x$.}
\end{enumerate}

\vspace{20pt}

\begin{warningbox}
Part (d) is the Gaussian density. The derivative you just computed is the
density times a linear function of $x$. This structure — density times score
— reappears in Module 02 when we differentiate the log-likelihood with respect
to a \emph{parameter}.
\end{warningbox}

\vspace{24pt}


% ============================================================
% PROBLEM 2 — Integration
% ============================================================
\begin{problembox}[title={Problem 2 \hfill $\star$}]
Evaluate each integral. Identify the technique you are using (direct
antiderivative, substitution, or integration by parts) before proceeding.
\end{problembox}

\begin{enumerate}
  \item $\displaystyle\int_0^1 x\, e^{-x^2/2}\, dx$
  \quad \textit{Use substitution.}

  \item $\displaystyle\int x \ln x\, dx$
  \quad \textit{Use integration by parts.}

  \item $\displaystyle\int_{-\infty}^{\infty} e^{-x^2/2}\, dx$
  \quad \textit{State the result; do not derive it. Explain why this integral
  is relevant to the Gaussian PDF.}

  \item $\displaystyle\int_0^{\infty} x\, e^{-x}\, dx$
  \quad \textit{Use integration by parts. Evaluate carefully at the boundaries.}
\end{enumerate}

\vspace{24pt}


% ============================================================
% PROBLEM 3 — Statistical Bridge: The Log-Likelihood
% ============================================================
\begin{problembox}[title={Problem 3 \hfill $\star\star$}]
Suppose we observe $n$ independent draws $x_1, x_2, \ldots, x_n$ from a
$\Normal(\mu, \sigma^2)$ distribution with \emph{known} variance $\sigma^2 = 1$.
The log-likelihood for the parameter $\mu$ is:

$$\loglik(\mu) = -\frac{n}{2}\ln(2\pi) - \frac{1}{2}\sum_{i=1}^n (x_i - \mu)^2$$
\end{problembox}

\begin{enumerate}
  \item Differentiate $\loglik(\mu)$ with respect to $\mu$.
  Show each step; treat $x_1, \ldots, x_n$ as constants.

  \item Set your derivative equal to zero and solve for $\hat{\mu}$.
  Express your answer in closed form.

  \item What is the name of the estimator you just derived?
  Does your answer depend on $\sigma^2$? Should it?

  \item Compute the second derivative $\dfrac{d^2\,\loglik}{d\mu^2}$.
  What does the sign tell you about whether $\hat{\mu}$ is a maximum or
  a minimum of $\loglik(\mu)$?
\end{enumerate}

\vspace{16pt}

\begin{bridgebox}
This is maximum likelihood estimation in its simplest form. In Module 02 you
will generalise exactly this calculation to $\vtheta \in \R^p$, replacing
the derivative with the gradient and the second derivative with the Hessian.
The structure — differentiate, set to zero, check the second-order condition
— is identical.
\end{bridgebox}

\vspace{24pt}


% ============================================================
% PROBLEM 4 — Notation Transcription (Extension)
% ============================================================
\begin{problembox}[title={Problem 4 \hfill $\star\star\star$ — Extension (Optional)}]
This problem requires no calculus. Its purpose is to build fluency with the
notation conventions established in Section~4 of the notes. Read
\texttt{assets/latex-preamble.tex} and the notation table in the notes, then
rewrite each expression in full English and identify the macro used to
typeset it.
\end{problembox}

\begin{enumerate}
  \item $\posterior$ \quad (identify what each term represents)

  \item $\MLE{\theta}$ \quad (what does the hat and subscript denote?)

  \item $\MVN{\vmu}{\mSigma}$ \quad (what are the two arguments, and what type
  of object is each?)

  \item Write the log-likelihood for a single observation $x$ from a
  $\Normal(\mu, \sigma^2)$ distribution using \emph{only} macros from the
  preamble (no raw \texttt{\textbackslash frac}, \texttt{\textbackslash exp},
  or \texttt{\textbackslash sigma} that are not wrapped in a course macro).
  Compare your expression with what appears in Problem~3.
\end{enumerate}

\vspace{24pt}


% ============================================================
% REFLECTION
% ============================================================
\begin{tcolorbox}[
  colback=ts-light,
  colframe=ts-green,
  title=Reflection,
  fonttitle=\bfseries\color{ts-green}
]
Before moving to the Python lab, take two minutes to answer these in your notes:

\begin{enumerate}
  \item Which differentiation rule required the most deliberate attention today?
  Describe the specific moment where you had to slow down.
  \item In Problem~3, you derived an estimator using only first-year calculus.
  What additional machinery would you need to do the same calculation for
  $\vtheta = (\mu_1, \mu_2, \ldots, \mu_p)\T$?
  \item Is there any macro in Section~4 of the notes whose meaning is still
  unclear? Write it down now so it surfaces during the Python lab.
\end{enumerate}

\medskip
\textbf{Next:} Python Lab (\texttt{python-lab.ipynb})
$\to$ Solutions (\texttt{solutions/module-00-orientation/unit-03-single-variable-review/})
\end{tcolorbox}


% ============================================================
\end{document}
% ============================================================
